\begin{abstract}
Modern software teams increasingly depend on two classes of assistive
technology to keep source code secure: static application security
testing (SAST) and, more recently, large language models (LLMs) used for
code review and vulnerability triage. Both are usually delivered as cloud
services, which obliges an organisation to transmit proprietary source
code to a third party and complicates adoption in regulated or air-gapped
environments. This paper presents \cs, a local-first desktop platform that
unifies pattern-based static analysis, containerised dynamic execution,
and on-device LLM reasoning into a single privacy-preserving workflow for
auditing the security and structural health of a source-code repository.
Every stage runs on the developer's machine: pattern detectors and a
cyclomatic-complexity engine operate in process, a quantised
\mbox{Llama~3.2} model performs semantic review inside a local container,
and a resource-limited Docker sandbox exercises the project at runtime.
Findings from the three modalities are fused, de-duplicated, and combined
with software metrics into a single \phindex. We describe the
architecture, the analysis pipeline, and the scoring model, and we
present the twelve-screen graphical interface through which a developer
drives and inspects an audit. A worked case study on a representative
retrieval-augmented-generation service ($47$ files, ${\sim}5.2$~KLOC)
shows the platform surfacing $20$ findings spanning credential handling,
prompt and query injection, unsafe deserialisation, and server-side
request forgery, alongside five cyclomatic hot-spots and a moderate
composite risk score, while completing end to end in under a minute and
without a single outbound network request. We discuss the trade-offs of
local inference, the limitations of the current prototype, and directions
for a larger empirical evaluation.
\end{abstract}

\begin{IEEEkeywords}
Static application security testing, large language models, on-device
inference, dynamic analysis, cyclomatic complexity,
retrieval-augmented generation, privacy-preserving software tools,
developer tooling.
\end{IEEEkeywords}
